% !TEX root = task-3-6-1.tex
% Editable task bank. Append new BankTask blocks at the end of the document.
\documentclass[12pt,a4paper]{article}
\usepackage{physics-textbook}
\usepackage{graphicx,needspace}
\usetikzlibrary{decorations.pathmorphing}
\newif\ifShowSolutions
\ifdefined\TutorEdition
  \ShowSolutionstrue
\else
  \ShowSolutionsfalse % Change to \ShowSolutionstrue to display solutions.
\fi
\newcounter{banktask}
\newsavebox{\BankTaskBox}
\newcommand{\BankTaskContent}[3]{%
  \noindent\textbf{\thebanktask.}\quad #2\par
  \ifShowSolutions
    \Needspace{50mm}% Keep the solution heading with its opening text or diagram.
    #3\par
    \ifdefined\TutorEdition
      \ifcsname TutorSource#1\endcsname
        \nopagebreak\smallskip{\small\csname TutorSource#1\endcsname\par}%
      \fi
    \fi
  \fi
}
\newcommand{\BankTask}[3]{%
  \refstepcounter{banktask}\par\medskip
  \sbox{\BankTaskBox}{\begin{minipage}{\linewidth}%
    \BankTaskContent{#1}{#2}{#3}\end{minipage}}%
  \ifdim\dimexpr\ht\BankTaskBox+\dp\BankTaskBox\relax>\dimexpr\textheight-20mm\relax
    \Needspace{35mm}\BankTaskContent{#1}{#2}{#3}%
  \else
    \Needspace{\dimexpr\ht\BankTaskBox+\dp\BankTaskBox\relax}%
    \noindent\usebox{\BankTaskBox}\par
  \fi
  \medskip
}
\begin{document}
\BookSubsection{3.6.1}{Диэлектрическая проницаемость}


% Source sdamgia-31415: https://phys-ege.sdamgia.ru/problem?id=31415
\BankTask{sdamgia-31415}{
Стеклянные незаряженные кубики~$1$ и $2$ и алюминиевые незаряженные кубики~$3$ и $4$ сблизили вплотную и поместили в электрическое поле, напряженность которого направлена горизонтально влево, как показано в левой части рисунка. Затем кубики раздвинули и уже после этого выключили электрическое поле (правая часть рисунка).\par
$1$. В электрическом поле (левая часть рисунка) часть электронов перешла из кубика~$4$ в кубик~$3$.\par
$2$. Кубик $2$, изображенный в правой части рисунка, не заряжен.\par
$3$. В электрическом поле (левая часть рисунка) правая грань кубика~$2$ заряжена положительно.\par
$4$. Кубик $3$, изображенный в правой части рисунка, имеет положительный заряд.\par
$5$. В электрическом поле (левая часть рисунка) разность потенциалов между левой гранью кубика~$3$ и правой гранью кубика~$4$ равна нулю.\par

\begin{center}\begin{tikzpicture}[book diagram,>=stealth]
\draw(0,0)rectangle(1.2,1.2);\draw(.6,0)--(.6,1.2)(0,.6)--(1.2,.6);\foreach\x/\y/\s in{.3/.9/1,.9/.9/2,.3/.3/3,.9/.3/4}\node at(\x,\y){$\s$};\foreach\y in{.3,.5,.7,.9}\draw[->](-.2,\y)--(-1,\y);\node[above]at(-.7,1){$\vec E$};\draw[dashed](1.65,-.2)--(1.65,1.4)(1.8,-.2)--(1.8,1.4);\foreach\x/\y/\s in{2.1/.8/1,3.1/.8/2,2.1/-.2/3,3.1/-.2/4}{\draw(\x,\y)rectangle++(.6,.6);\node at(\x+.3,\y+.3){$\s$};}
\end{tikzpicture}\end{center}
\par
}{
\smallskip\noindent\textbf{Решение}\par
$1$. Неверно. Электроны под действием электрического поля движутся против силовых линий, поэтому переместятся вправо, тогда часть электронов перейдет с кубика~$3$ в кубик~$4$.\par
$2$. Верно. Стекло является диэлектриком, в котором нет свободных носителей заряда. Поэтому кубики из стекла останутся незаряженными.\par
$3$. Неверно. В кубиках из стекла произойдет поляризация, в результате которой правая часть кубика~$2$ будет иметь отрицательный заряд.\par
$4$. Верно. Поскольку электроны в металлах переместятся против силовых линий, то в кубике~$3$ будет недостаток электронов, он будет заряжен положительно.\par
$5$. Верно. В проводнике, помещенном в электростатическое поле, напряженность станет равной $0$, поэтому и разность потенциалов между левой гранью кубика~$3$ и правой гранью кубика~$4$ равна нулю.\par
Ответ: $245$.\par
}

% Source sdamgia-3515: https://phys-ege.sdamgia.ru/problem?id=3515
\BankTask{sdamgia-3515}{
Между пластинами заряженного плоского конденсатора поместили диэлектрик с диэлектрической проницаемостью $\varepsilon$ так, что он полностью заполнил объем между пластинами. Как изменились емкость конденсатора, заряд на пластинах и напряжение между ними, если конденсатор отключен от источника?\par

\begin{center}\begin{minipage}[t]{.47\linewidth}\textbf{Физическая характеристика}\par\smallskip А) Заряд на пластинах\par
Б) Напряжение между пластинами\par
В) Емкость конденсатора\end{minipage}\hfill\begin{minipage}[t]{.47\linewidth}\textbf{Ее изменение}\par\smallskip $1$) Уменьшится в $\varepsilon$ раз\par
$2$) Останется неизменной\par
$3$) Увеличится в $\varepsilon$ раз\end{minipage}\end{center}
\par
\begin{center}\small\setlength{\tabcolsep}{3pt}\begin{tabular}{|p{0.3\linewidth}|p{0.3\linewidth}|p{0.3\linewidth}|}\hline
$A$ & Б & В\\\hline
 &  & \\\hline
\end{tabular}\end{center}\par
}{
\smallskip\noindent\textbf{Решение}\par
Поскольку конденсатор отключен от источника, заряд на его пластинах меняться не может, ему просто некуда деться и неоткуда взяться (А — $2$). Емкость плоского конденсатора пропорциональна диэлектрической проницаемости диэлектрика внутри него: $C= \frac{ \varepsilon\varepsilon_{0}S}{ d } .$ Следовательно, введение диэлектрика приводит к увеличению емкости конденсатора в $\varepsilon$ раз (В — $3$). Наконец, напряжение между пластинами конденсатора связано с его емкостью и зарядом конденсатора соотношением: $U= \frac{ q}{ C } .$ Таким образом, после введения диэлектрика напряжением уменьшится в $\varepsilon$ раз (Б — $1$).\par
Ответ: $213$.\par
}

% Source sdamgia-3516: https://phys-ege.sdamgia.ru/problem?id=3516
\BankTask{sdamgia-3516}{
Между пластинами заряженного плоского конденсатора поместили диэлектрик с диэлектрической проницаемостью $\varepsilon$ так, что он полностью заполнил объем между пластинами. Как изменились емкость конденсатора, заряд на пластинах и напряжение между ними, если конденсатор подключен к источнику?\par

\begin{center}\begin{minipage}[t]{.47\linewidth}\textbf{Физическая характеристика}\par\smallskip А) Заряд на пластинах\par
Б) Напряжение между пластинами\par
В) Емкость конденсатора\end{minipage}\hfill\begin{minipage}[t]{.47\linewidth}\textbf{Ее изменение}\par\smallskip $1$) Уменьшится в $\varepsilon$ раз\par
$2$) Останется неизменной\par
$3$) Увеличится в $\varepsilon$ раз\end{minipage}\end{center}
\par
\begin{center}\small\setlength{\tabcolsep}{3pt}\begin{tabular}{|p{0.3\linewidth}|p{0.3\linewidth}|p{0.3\linewidth}|}\hline
$A$ & Б & В\\\hline
 &  & \\\hline
\end{tabular}\end{center}\par
}{
\smallskip\noindent\textbf{Решение}\par
Поскольку конденсатор подключен к источнику, напряжение между пластинами конденсатора равно напряжению источника, а значит, постоянно (Б — $2$). Емкость плоского конденсатора пропорциональна диэлектрической проницаемости диэлектрика внутри него: $C= \frac{ \varepsilon\varepsilon_{0}S}{ d } .$ Следовательно, введение диэлектрика приводит к увеличению емкости конденсатора в $\varepsilon$ раз (В — $3$). Наконец, заряд на пластинах конденсатора связан с его емкостью и напряжением на нем соотношением: $q=CU.$ Таким образом, после введения диэлектрика заряд конденсатора увеличится в $\varepsilon$ раз (А — $3$).\par
Ответ: $323$.\par
}

% Source sdamgia-10474: https://phys-ege.sdamgia.ru/problem?id=10474
\BankTask{sdamgia-10474}{
Две параллельные металлические пластины больших размеров расположены на расстоянии $d$ друг от друга и подключены к источнику постоянного напряжения (рис. $1$). Пластины закрепили на изолирующих подставках и спустя длительное время отключили от источника (рис. $2$).\par
$1$. Напряженность электрического поля в точке А больше, чем в точке В.\par
$2$. Потенциал электрического поля в точке А больше, чем в точке С.\par
$3$. Если увеличить расстояние между пластинами $d$, то напряженность электрического поля в точке С не изменится.\par
$4$. Если уменьшить расстояние между пластинами $d$, то заряд правой пластины не изменится.\par
$5$. Если пластины полностью погрузить в керосин, то энергия электрического поля конденсатора останется неизменной.\par

\begin{center}\begin{tikzpicture}[book diagram,>=stealth]
\foreach\shift in{0,4}{\begin{scope}[xshift=\shift cm]\draw[thick](0,0)--(0,2.6)(2,0)--(2,2.6);\draw[<->](0,3)--(2,3)node[midway,above]{$d$};\draw(0,2.6)--(0,3.2)(2,2.6)--(2,3.2);\end{scope}}\draw(0,1.3)--(-.8,1.3)--(-.8,-.8)--(.65,-.8);\draw(.65,-.8)circle(.1);\draw(1.35,-.8)circle(.1);\draw(1.35,-.8)--(2.8,-.8)--(2.8,1.3)--(2,1.3);\node[above]at(.65,-.7){$-$};\node[above]at(1.35,-.7){$+$};\node at(1,-1.3){Рис.\ 1};\draw(4,1.3)--(3.5,1.3)(6,1.3)--(6.5,1.3);\foreach\x/\y/\s in{4.55/.55/A,4.55/1.7/B,5.3/1.15/C}{\fill(\x,\y)circle(1.5pt);\node[above]at(\x,\y){$\s$};}\node at(5,-1.3){Рис.\ 2};
\end{tikzpicture}\end{center}
\par
}{
\smallskip\noindent\textbf{Решение}\par
После того как длительное время пластины были подключены к источнику постоянного напряжения, они зарядились: левая пластина — отрицательно, правая — положительно.\par
$1$. Внутри плоского заряженного конденсатора электрическое поле однородно. Напряженности поля в точках А и В одинаковые. Утверждение $1$ неверно.\par
$2$. Потенциал электрического поля внутри конденсатора убывает от положительной пластины к отрицательной. Потенциал электрического поля в точке А меньше, чем в точке С. Утверждение $2$ неверно.\par
$3$. Поскольку пластины отключены от источника, то заряд и его поверхностная плотность на них не меняется при изменении расстояния. Значит, не будет изменяться и напряженность электрического поля между пластинами. Утверждение $3$ верно.\par
$4$. Заряд пластин остается постоянным, независимо от того, сдвигают пластины или нет. Утверждение $4$ верно.\par
$5$. Диэлектрическая проницаемость керосина больше $1$. При полном погружении в керосин энергия электрического поля конденсатора $W= \frac{ q ^2 }{ 2C }$ уменьшится. Утверждение $5$ неверно.\par
Ответ: $34$.\par
}

% Source sdamgia-24107: https://phys-ege.sdamgia.ru/problem?id=24107
\BankTask{sdamgia-24107}{
Плоский конденсатор состоит из двух квадратных пластин. Пространство между пластинами заполнено однородным диэлектриком. Установите соответствие между зависимостями емкости $C$ этого конденсатора от физических величин и графиками, изображающими эти зависимости.\par

\begin{center}\begin{minipage}[t]{.47\linewidth}\textbf{Зависимость}\par\smallskip А) емкости С конденсатора от длины $a$ стороны пластины\par
Б) емкости С конденсатора от диэлектрической проницаемости $\varepsilon$ диэлектрика\end{minipage}\hfill\begin{minipage}[t]{.47\linewidth}\textbf{График зависимости}\par\smallskip 
\begin{center}\begin{tikzpicture}[book diagram,scale=.65,transform shape,>=stealth]
\begin{scope}[shift={(0,3.4)}]\draw[task grid,xstep=.6,ystep=.2](0,0)grid(2.4,2.4);\draw[task axis](0,0)--(2.7,0);\draw[task axis](0,0)--(0,2.7)node[left]{$C$};\node[below left]at(0,0){$0$};\node[left]at(-.2,2.4){$ 1) $};\draw[task curve](0,0)--(1.8,2.4);\end{scope}\begin{scope}[shift={(3.6,3.4)}]\draw[task grid,xstep=.6,ystep=.2](0,0)grid(2.4,2.4);\draw[task axis](0,0)--(2.7,0);\draw[task axis](0,0)--(0,2.7)node[left]{$C$};\node[below left]at(0,0){$0$};\node[left]at(-.2,2.4){$ 2) $};\draw[task curve](.6,.8)--(2.4,2.4);\end{scope}\begin{scope}[shift={(0,0)}]\draw[task grid,xstep=.6,ystep=.2](0,0)grid(2.4,2.4);\draw[task axis](0,0)--(2.7,0);\draw[task axis](0,0)--(0,2.7)node[left]{$C$};\node[below left]at(0,0){$0$};\node[left]at(-.2,2.4){$ 3) $};\draw[task curve,domain=0:2.08,samples=60]plot(\x,{.5556*\x*\x});\end{scope}\begin{scope}[shift={(3.6,0)}]\draw[task grid,xstep=.6,ystep=.2](0,0)grid(2.4,2.4);\draw[task axis](0,0)--(2.7,0);\draw[task axis](0,0)--(0,2.7)node[left]{$C$};\node[below left]at(0,0){$0$};\node[left]at(-.2,2.4){$ 4) $};\draw[task curve,domain=.22:2.4,samples=60]plot(\x,{.52/\x});\end{scope}
\end{tikzpicture}\end{center}
\end{minipage}\end{center}
\par
\begin{center}\small\setlength{\tabcolsep}{3pt}\begin{tabular}{|p{0.45\linewidth}|p{0.45\linewidth}|}\hline
А & Б\\\hline
 & \\\hline
\end{tabular}\end{center}\par
}{
\smallskip\noindent\textbf{Решение}\par
Емкость плоского конденсатора с квадратными пластинами равна $C= \frac{ \varepsilon\varepsilon_{0}S}{ d } = \frac{ \varepsilon\varepsilon_{0}a ^2 }{ d } .$ Значит, емкость пропорциональна квадрату стороны пластины. Графиком квадратичной зависимости является график $3$.\par
У воздушного конденсатора диэлектрическая проницаемость равна $1$. При увеличении диэлектрической проницаемости емкость конденсатора увеличивается прямо пропорционально. Меньше $1$ быть не может. Учитывая выше сказанное, графиком зависимости емкости конденсатора от диэлектрической проницаемости является график $2$.\par
Ответ: $32$.\par
}

% Source sdamgia-9067: https://phys-ege.sdamgia.ru/problem?id=9067
\BankTask{sdamgia-9067}{
Необходимо экспериментально изучить зависимость емкости плоского конденсатора от свойств диэлектрика, помещенного между его пластинами. На всех представленных ниже рисунках $S$ — площадь пластины конденсатора, $d$ — расстояние между пластинами. Какие две установки следует использовать для проведения такого исследования?\par

\begin{center}\begin{tikzpicture}[book diagram,>=stealth]
\node[draw,circle]at(0,1.5){$ 1 $};\fill[orange!25](-0.6,-0.9)rectangle(0.6,0.9);\draw(-0.6,-0.9)--(-0.6,0.9)(0.6,-0.9)--(0.6,0.9);\draw(-0.6-.7,0)--(-0.6,0)(0.6,0)--(0.6+.7,0);\node[above]at(-0.6-.4,0){$ S $};\node[rotate=90]at(0,0){стекло};\draw[<->](-0.6,-0.9-.5)--(0.6,-0.9-.5)node[midway,above]{$ d $};
\end{tikzpicture}\end{center}
\par

\begin{center}\begin{tikzpicture}[book diagram,>=stealth]
\node[draw,circle]at(0,1.5){$ 2 $};\fill[orange!25](-1.2,-0.9)rectangle(1.2,0.9);\draw(-1.2,-0.9)--(-1.2,0.9)(1.2,-0.9)--(1.2,0.9);\draw(-1.2-.7,0)--(-1.2,0)(1.2,0)--(1.2+.7,0);\node[above]at(-1.2-.4,0){$ S $};\node[rotate=0]at(0,0){керамика};\draw[<->](-1.2,-0.9-.5)--(1.2,-0.9-.5)node[midway,above]{$ 2d $};
\end{tikzpicture}\end{center}
\par

\begin{center}\begin{tikzpicture}[book diagram,>=stealth]
\node[draw,circle]at(0,1.5){$ 3 $};\fill[orange!25](-0.6,-0.9)rectangle(0.6,0.9);\draw(-0.6,-0.9)--(-0.6,0.9)(0.6,-0.9)--(0.6,0.9);\draw(-0.6-.7,0)--(-0.6,0)(0.6,0)--(0.6+.7,0);\node[above]at(-0.6-.4,0){$ S $};\node[rotate=90]at(0,0){керамика};\draw[<->](-0.6,-0.9-.5)--(0.6,-0.9-.5)node[midway,above]{$ d $};
\end{tikzpicture}\end{center}
\par

\begin{center}\begin{tikzpicture}[book diagram,>=stealth]
\node[draw,circle]at(0,2.4){$ 4 $};\fill[orange!25](-0.6,-1.8)rectangle(0.6,1.8);\draw(-0.6,-1.8)--(-0.6,1.8)(0.6,-1.8)--(0.6,1.8);\draw(-0.6-.7,0)--(-0.6,0)(0.6,0)--(0.6+.7,0);\node[above]at(-0.6-.4,0){$ 2S $};\node[rotate=90]at(0,0){стекло};\draw[<->](-0.6,-1.8-.5)--(0.6,-1.8-.5)node[midway,above]{$ d $};
\end{tikzpicture}\end{center}
\par

\begin{center}\begin{tikzpicture}[book diagram,>=stealth]
\node[draw,circle]at(0,2.4){$ 5 $};\fill[orange!25](-1.2,-1.8)rectangle(1.2,1.8);\draw(-1.2,-1.8)--(-1.2,1.8)(1.2,-1.8)--(1.2,1.8);\draw(-1.2-.7,0)--(-1.2,0)(1.2,0)--(1.2+.7,0);\node[above]at(-1.2-.4,0){$ 2S $};\node[rotate=0]at(0,0){керамика};\draw[<->](-1.2,-1.8-.5)--(1.2,-1.8-.5)node[midway,above]{$ 2d $};
\end{tikzpicture}\end{center}
\par
}{
\smallskip\noindent\textbf{Решение}\par
Емкость конденсатора прямо пропорциональна площади его пластин, диэлектрической проницаемости и обратно пропорциональна расстоянию между обкладками: $C= \frac{ \varepsilon\varepsilon _{0}S}{ d } .$\par
Для экспериментального изучения емкости конденсатора в зависимости от свойств диэлектрика необходимо выбрать два конденсатора с разными диэлектриками между обкладками, но с одинаковыми остальными параметрами. Для этого подходят конденсаторы под номерами $1$ и $3$.\par
Ответ: $13$.\par
}

% Source sdamgia-4360: https://phys-ege.sdamgia.ru/problem?id=4360
\BankTask{sdamgia-4360}{
Плоский заряженный воздушный конденсатор, отключенный от источника напряжения, заполняют диэлектриком. Какова диэлектрическая проницаемость диэлектрика, если напряженность электрического поля в диэлектрике между пластинами заполненного конденсатора меньше напряженности электрического поля незаполненного конденсатора в $1{,}25$ раза?\par
}{
\smallskip\noindent\textbf{Решение}\par
Поскольку конденсатор отключен от источника заряд его пластин остается неизменным. Емкость конденсатора в каждом случае равна $C_{1}= \frac{ \varepsilon _{0}S}{ d }$ и $C_{2}= \frac{ \varepsilon \varepsilon _{0}S}{ d } .$ Напряжение между пластинами конденсатора в каждом случае $U_{1}= \frac{ q}{ C_{1} }$ и $U_{2}= \frac{ q}{ C_{2} } .$\par
Напряженность поля внутри конденсатора $E_{1}= \frac{ U_{1}}{ d } = \frac{ q}{ \varepsilon _{0}S }$ и $E_{2}= \frac{ U_{2}}{ d } = \frac{ q}{ \varepsilon \varepsilon _{0}S } .$ По условию задачи $E_{1}=1{,}25 E_{2}.$ Следовательно, диэлектрическая проницаемость вещества равна $\varepsilon =1{,}25.$\par
Ответ: $1{,}25$.\par
}

% Source sdamgia-38859: https://phys-ege.sdamgia.ru/problem?id=38859
\BankTask{sdamgia-38859}{
Плоский воздушный вертикальный конденсатор подключен к источнику с постоянным напряжением. Не отключая конденсатор от источника его погружают в жидкий диэлектрик с диэлектрической проницаемостью $\varepsilon=4,$ при этом диэлектрик заполняет $\frac{ 1}{ 3 }$ объема конденсатора. Во сколько раз изменилась энергия электрического поля конденсатора после погружения его в жидкий диэлектрик?\par
}{
\smallskip\noindent\textbf{Решение}\par
Первоначально конденсатор имел емкость $C_{0}= \frac{ \varepsilon _{0}S}{ d } ,$ где $S$ — площадь одной пластины, $d$ — расстояние между пластинами. В этом случае энергия конденсатора $W_{0}= \frac{ C_{0}U ^2 }{ 2 } .$\par
После погружения части конденсатора в жидкий диэлектрик заполненная часть конденсатора по условию $V_{1}= \frac{ 1}{ 3 } V_{0}= \frac{ 1}{ 3 } Sd,$ а незаполненная $V_{2}= \frac{ 2}{ 3 } V_{0}= \frac{ 2}{ 3 } Sd.$ Учитывая, что расстояние между пластинами не изменяется, получаем соотношения: $S_{1}= \frac{ 1}{ 3 } S$ и $S_{2}= \frac{ 2}{ 3 } S.$\par
Таким образом, получили два конденсатора, емкости которых:\par
\[C_{1}= \frac{ \varepsilon \varepsilon _{0}S}{ 3d } = \frac{ \varepsilon C_{0}}{ 3 } ,\]
\[C_{2}= \frac{ 2 \varepsilon _{0}S}{ 3d } = \frac{ 2C_{0}}{ 3 } .\]
При этом получившиеся конденсаторы соединены параллельно, поэтому общая емкость получившихся конденсаторов:\par
\[C=C_{1} + C_{2}= \frac{ \varepsilon C_{0}}{ 3 } + \frac{ 2C_{0}}{ 3 } = \frac{ ( \varepsilon + 2 ) C_{0}}{ 3 } .\]
Поскольку конденсатор не отключали от источника тока, то напряжение на нем не изменилось. Тогда энергия заряженных конденсаторов после погружения в диэлектрик стала равной:\par
\[W= \frac{ CU ^2 }{ 2 } = \frac{ ( \varepsilon + 2 ) C_{0}U ^2 }{ 6 } .\]
Найдем их отношение:\par
\[\frac{ W}{ W_{0} } = \frac{ 2 + \varepsilon }{ 3 } = \frac{ 2 + 4}{ 3 } =2.\]
Энергия заряженного конденсатора увеличилась в $2$ раза.\par
Ответ: увеличилась в $2$ раза.\par
}

% Source sdamgia-6434: https://phys-ege.sdamgia.ru/problem?id=6434
\BankTask{sdamgia-6434}{
Два плоских конденсатора емкостью $C$ и $2C$ соединили параллельно и зарядили до напряжения $U$. Затем ключ К разомкнули, отключив конденсаторы от источника (см. рис.). Пространство между их обкладками заполнено жидким диэлектриком с диэлектрической проницаемостью $\varepsilon$. Какой будет разность потенциалов между обкладками, если из левого конденсатора диэлектрик вытечет?\par

\begin{center}\begin{tikzpicture}[book diagram,>=stealth]
\draw(0,0)--(0,1.1)(0,1.3)--(0,2)--(3,2)--(3,1.3)(3,1.1)--(3,0)--(0,0);\draw[fill=orange!25](-.3,1.1)rectangle(.3,1.3);\draw[fill=orange!25](2.7,1.1)rectangle(3.3,1.3);\node[left]at(0,1.6){$C$};\node[right]at(3,1.6){$2C$};\node[left]at(0,.8){$\varepsilon$};\node[right]at(3,.8){$\varepsilon$};\draw(1.5,2)--(1.5,2.5)--(5,2.5)--(5,1.8);\draw[very thick](4.6,1.8)--(5.4,1.8)(4.8,1.6)--(5.2,1.6);\node[right]at(5.35,1.9){$\mathcal E$};\draw(5,1.6)--(5,.8);\draw(5,.1)--(5,0)--(1.5,0);\draw[fill=white](5,.8)circle(2pt)(5,.1)circle(2pt);\draw(5,.1)--(5.45,.6);\node[right]at(5.35,.35){$K$};
\end{tikzpicture}\end{center}
\par
}{
\smallskip\noindent\textbf{Решение}\par
Заряд на конденсаторах равен $q=C_{\text{общ}}U,$ при параллельном соединении конденсаторов их общая емкость равна $C_{\text{общ}}=C + 2C=3C.$\par
\[q=3CU, ( 1 )\]
где $3C$ — суммарная емкость конденсаторов, когда оба они заполнены жидким диэлектриком. После вытекания диэлектрика из левого конденсатора суммарный заряд останется прежним. Емкость плоского конденсатора равна $C= \frac{ \varepsilon \varepsilon _{0}S}{ d } ,$ поэтому емкость левого конденсатора станет равной $\frac{ C}{ \varepsilon } ,$ а суммарная емкость станет равна $C/\varepsilon+2C$, а напряжение будет равно $U_{1}$, так что:\par
\[q= ( \frac{ C}{ \varepsilon } + 2C ) U_{1}. ( 2 )\]
Решая систему уравнений ($1$) и ($2$), получим ответ:\par
\[U_{1}= \frac{ 3\varepsilon}{ 1 + 2\varepsilon } U.\]
Ответ: $U_{1}= \frac{ 3\varepsilon}{ 1 + 2\varepsilon } U.$\par
}

% Source sdamgia-6662: https://phys-ege.sdamgia.ru/problem?id=6662
\BankTask{sdamgia-6662}{
В плоский воздушный конденсатор емкостью $16 \,\text{мкФ}$ вводят пластину с диэлектрической проницаемостью, равной $4$, после чего заряжают конденсатор, подключив его к клеммам источника с напряжением $6 \,\text{В}$. На сколько уменьшится энергия этого конденсатора, если, не отсоединяя конденсатор от источника, извлечь пластину из конденсатора? (В микроджоулях.)\par
}{
\smallskip\noindent\textbf{Решение}\par
Емкость плоского конденсатора вычисляется по формуле: $C= \frac{ \varepsilon_{0}\varepsilon S}{ d } ,$ где $\varepsilon_{0}$ — электрическая постоянная, $\varepsilon$ — диэлектрическая проницаемость материала внутри конденсатора, $S$ — площадь конденсатора, $d$ — расстояние между пластинами конденсатора. Диэлектрическая проницаемость воздуха равна единице. Значит, при введении пластины с диэлектрической проницаемостью $4$ емкость конденсатора увеличится в четыре раза, то есть станет равной $64 \,\text{мкФ}$.\par
Энергия конденсатора в этом случае будет равна: $W_{1}= \frac{ C_{1}U ^2 }{ 2 } = \frac{ 64\,\text{мкФ} \cdot ( 6\,\text{В} ) ^2 }{ 2 } =1152\,\text{мкДж}.$ Конденсатор не отключают от источника, поэтому при извлечении пластины из конденсатора напряжение на обкладках конденсатора остается прежним. Тогда энергия конденсатора будет равна: $W_{2}= \frac{ C_{2}U ^2 }{ 2 } = \frac{ 16\,\text{мкФ} \cdot ( 6\,\text{В} ) ^2 }{ 2 } =288\,\text{мкДж}.$\par
Разность этих энергий: $\Delta W=1152\,\text{мкДж} - 288\,\text{мкДж} = 864\,\text{мкДж}.$\par
Ответ: $864 \,\text{мкДж}$.\par
}

\end{document}
